\documentclass[reqno,12pt]{amsproc}

\usepackage{amsmath}
%\documentclass[reqno,11pt]{amsart}
%\usepackage[a4paper,hypertex]{hyperref}
%\AtBeginDocument{\noindent\small
%http://www.ma.utexas.edu/mpmarc/mp_arc-home.html
%\documentclass[reqno]{amsart}
%mmmmmmmmmmmmmmmmmmmmmmmmmmmmmmmmmmmmmmmmmmmmmmmmmmmmmmmmmmmmmmmmmmmmmmm
\newtheorem{theo}{Theorem}
\newtheorem{rem}{Remark}
\newtheorem{ex}{Example}
\newtheorem{lem}{Lemma}
\newtheorem{conj}{Conjecture}
\newtheorem{prop}{Proposition}
\newtheorem{cor}{Corollary}
\newtheorem{df}{Definition}
\newtheorem{hyp}{Hypothesis}
\newtheorem{ass}{Assumption}
%mmmmmmmmmmmmmmmmmmmmmmmmmmmmmmmmmmmmmmmmmmmmmmmmmmmmmmmmmmmmmmmmmmmmmmm
\newcommand\eps\varepsilon
\newcommand\ph\varphi
\newcommand\kap\varkappa
\newcommand\R {\mathbb{R}}
\newcommand\T {\mathbb{T}}
\newcommand\Z {\mathbb{Z}}
\newcommand\N {\mathbb{N}}
\newcommand\bC {\mathbb{C}}
\newcommand\bs {\boldsymbol}
%\hypersetup{colorlinks=true, urlcolor=blue, linkcolor=blue, citecolor=blue, breaklinks=true}

\newcommand{\bigsize}{\fontsize{18pt}{20pt}\selectfont}
\usepackage{hyperref}
\usepackage{url}

\renewcommand{\refname}{References}
\renewcommand{\Re}{\mbox{\rm Re}\,}
\renewcommand{\Im}{\mbox{\rm Im}\,}
\usepackage{enumerate}
\usepackage{graphicx}
\usepackage{tikz-cd}
\usepackage{float}
\usepackage{yfonts}
\usepackage{mathrsfs}
\usepackage{amsthm}
\usepackage{amscd}




%mmmmmmmmmmmmmmmmmmmmmmm
\hoffset -35mm
\voffset 0mm
\overfullrule 0pt
\tolerance 1000

\textwidth=510pt
\textheight=650pt
%mmmmmmmmmmmmmmmmmmmmmmm

\begin{document}\title
{The D'Alembert-Lagrange Principle and Lagrange Equations}
\author[Oleg Zubelevich]{Oleg Zubelevich\\Faculty of Mechanics and Mathematics,\\
Lomonosov Moscow State University,\\ 1 Leninskie Gory, Moscow 119991, Russia}
\email{oezubel@gmail.com}

\date{}
%\thanks{The research was funded by a grant from the Russian Science
%Foundation (Project No. 19-71-30012)}
%\subjclass[2000]{70-01, 70H03, 70H30}
\subjclass[2020]{Primary 70-01; Secondary 70H03, 70H30}
\keywords{Classical mechanics, variational derivative, ideal constraints}


\begin{abstract}This expository article serves as a methodical guide, presenting a rigorous mathematical formulation of the D'Alembert-Lagrange principle and the derivation of the Lagrange equations for mechanical systems subject to ideal constraints. Designed for educational purposes, the paper establishes a clear geometric framework for the extended phase space, defining the space of virtual displacements and constraint reactions independently of their specific analytical representations. A central focus of this instructional text is the transition from the general equation of dynamics to the Lagrange equations of the second kind for holonomic systems. Specifically, we demonstrate that the Lagrange equations can be naturally and elegantly derived from considerations of the covariance of the variational derivative under the embedding mapping of the configuration manifold.
\end{abstract}

\maketitle

%%%%%%%%%%%  июнь  2016   %%%%%%%%%%%%%%%%%%%%%%
\section{Introduction}

Classical mechanics provides a profound framework for understanding the dynamics of physical systems. At the heart of analytical mechanics lie the D'Alembert-Lagrange principle and the Lagrange equations, which offer a powerful alternative to the direct application of Newton's laws, particularly when dealing with constrained systems.

Traditionally, in physics literature, these concepts are often introduced relying heavily on physical intuition, coordinate-dependent representations, and infinitesimal virtual displacements. However, a modern mathematical treatment demands a more rigorous, geometric perspective.

In this paper, we revisit the foundational principles of constrained mechanical systems to provide a mathematically precise exposition.

\section{Description of the problem at a physical level of rigor}

Let a system of point masses with masses
$m_1,\ldots,m_\nu$ be given.
Let $\boldsymbol r_1,\ldots,\boldsymbol r_\nu$ denote the radius vectors of these points in some inertial Cartesian system $OXYZ$:
$$\boldsymbol r_i=(X_i,Y_i,Z_i).$$
We represent the force acting on the $i$-th point as the sum
$\boldsymbol F_i+\boldsymbol R_i,$ where
\begin{align}\boldsymbol F_i&=\boldsymbol F_i(t,\boldsymbol r_1,\ldots,\boldsymbol r_{\nu},\boldsymbol {\dot r}_1,\ldots,\boldsymbol {\dot r}_{\nu})=(F_i^x,F_i^y,F_i^z),\nonumber\\
 \boldsymbol R_i&=\boldsymbol R_i(t,\boldsymbol r_1,\ldots,\boldsymbol r_{\nu},\boldsymbol {\dot r}_1,\ldots,\boldsymbol {\dot r}_{\nu})=(R_i^x,R_i^y,R_i^z).\nonumber
 \end{align}
The forces $\boldsymbol F_i$ are assumed to be given and are called active forces. 

The forces $\boldsymbol R_i$ compel the system of point masses to move in accordance with the equations
$$\Phi^k(t,\boldsymbol r_1,\ldots,\boldsymbol r_{\nu},\boldsymbol {\dot r}_1,\ldots,\boldsymbol {\dot r}_{\nu})=0,\quad k=1,\ldots, n,$$ which are called constraint equations or constraints. The forces $\boldsymbol R_i$ are called forces of constraint reactions or simply reactions. Reactions are to be determined during the solution of the problem. If there are no constraints, the reaction forces are set to zero.

Let us write Newton's second law for each point:
\begin{equation}\label{dfgjjj}m_i\boldsymbol {\ddot r}_i=\boldsymbol F_i+\boldsymbol R_i,\quad i=1,\ldots, {\nu}.\end{equation}
To transition to rigorous formulations, we introduce the following notations:
\begin{align}G&=\mathrm{diag}\,(m_1,m_1,m_1,\ldots,m_{\nu},m_{\nu},m_{\nu});\quad x=(X_1,Y_1,Z_1,\ldots,X_{\nu},Y_{\nu},Z_{\nu})^T;\nonumber\\
f(t,x,\dot x)&=(F_1^x,F_1^y,F_1^z,\ldots ,F_{\nu}^x,F_{\nu}^y,F_{\nu}^z),\quad N(t,x,\dot x)=(R_1^x,R_1^y,R_1^z,\ldots ,R_{\nu}^x,R_{\nu}^y,R_{\nu}^z),\nonumber\\
\ph^k(t,x,\dot x)&=\Phi^k.\nonumber
\end{align}



\section{The D'Alembert-Lagrange Principle}
The row vector $f$ and column vector $\ph=(\ph^1,\ldots,\ph^n)^T$ are assumed to be smooth functions\footnote{Hereafter, we assume all functions are as smooth as necessary for the formulas to make sense.} on the direct product $I\times D\times \mathbb{R}^m.$ Here $D\subset\mathbb{R}^m$ is a domain, $$ m=3\nu,\quad x=(x^1,\ldots, x^m)^T\in D,\quad t\in I=(t_1,t_2),\quad \dot x\in\mathbb{R}^m.$$
System (\ref{dfgjjj}) is written as follows:
\begin{equation}\label{seg44}
G\ddot x=f^T(t,x,\dot x)+N^T(t,x,\dot x),\end{equation}
and the constraints take the form:
\begin{equation}\label{xfg55}
\ph(t,x,\dot x)=0.\end{equation}
The domain $D$ is called the configuration space; the domain $D\times\mathbb{R}^m\ni (x,\dot x)$ is called the phase space; the domain $I\times D\times\mathbb{R}^m\ni (t,x,\dot x)$ is called the extended phase space of system (\ref{seg44}).



We represent system (\ref{seg44}) in covariant form:
\begin{equation}\label{azzxxx}[\mathscr T]=f+N,\quad[\mathscr T]= \frac{d}{dt}\frac{\partial\mathscr{T}}{\partial \dot x}-\frac{\partial\mathscr{T}}{\partial  x}, \end{equation}
where
$$\mathscr T=\frac{1}{2}\sum_{i=1}^\nu m_i|\boldsymbol {\dot r}_i|^2=\frac{1}{2}\dot x^TG\dot x$$
is the kinetic energy of the system.
\begin{rem}The representation of system (\ref{seg44}) in the form (\ref{azzxxx})
plays a key role in the following exposition (see section \ref{cbbbbb}). Therefore, we will adhere to this form, although at this stage it may seem strange.\end{rem}



If there exists a function $\mathscr V(t,x,\dot x)$ such that
$f=[\mathscr{V}],$
then the active forces are called generalized potential forces.
The function $\mathscr V(t,x,\dot x)$ is called a generalized potential.
If the generalized potential does not depend on $\dot x$, it is called a potential or potential energy.



Assume that the constraints are non-degenerate:
$$ \mathrm{rang}\,\ph_{\dot x}=n<m,\quad\ph_{\dot x}=\begin{pmatrix}
\frac{\partial \ph^1}{\partial \dot x^1}  & \cdots & \frac{\partial \ph^1}{\partial \dot x^m} \\
\vdots    & \ddots & \vdots  \\
\frac{\partial \ph^n}{\partial \dot x^1} & \cdots & \frac{\partial \ph^n}{\partial \dot x^m}
\end{pmatrix}.$$
The set of solutions to equation (\ref{xfg55})
$$S=\{(t,x,\dot x)\mid\ph(t,x,\dot x)=0\},\quad \dim S=2m+1-n,$$ which we assume to be non-empty, is a smooth submanifold of the extended phase space.


\begin{theo}[Axiom of constraints]\label{as11}
There exists a unique row vector $N=N(t,x,\dot x)$ smooth on $I\times D\times \mathbb{R}^m$ such that

1) the functions $\ph$ are first integrals of system (\ref{azzxxx});

2) for any column vector $\xi\in\mathbb{R}^m$ such that
\begin{equation}\label{sdg-0}\ph_{\dot x}(t,x,\dot x)\xi=0\end{equation}
the equality
\begin{equation}\label{sfg43}N(t,x,\dot x)\xi =0\end{equation}
holds.
\end{theo}\begin{df}The linear space of vectors $\xi\in\mathbb{R}^m$ satisfying (\ref{sdg-0}), i.e., $\ker\ph_{\dot x}(t,x,\dot x)$, is called the space of virtual displacements. This space is associated with the values of the variables $(t,x,\dot x)$.


The dimension of the space of virtual displacements is called the number of degrees of freedom of the system
(\ref{azzxxx}), (\ref{xfg55}).
\end{df} Obviously, the number of degrees of freedom is equal to $m-n=:r$.



\begin{df} If the forces $N$ are chosen in accordance with Theorem \ref{as11}, then the system (\ref{azzxxx}), (\ref{xfg55}) is called a system with ideal constraints. The constraints (\ref{xfg55}) are called ideal, and the forces $N$ are called reactions of ideal constraints.\end{df}


Below, only systems with ideal constraints are discussed.

By Theorem \ref{as11}, the manifold $S$ is an invariant manifold of system (\ref{azzxxx}).

\subsubsection*{ Proof of Theorem \ref{as11}.} \begin{lem}\label{dfh6ll} Let $X,Y,Z$ be vector spaces and
$$A:X\to Y,\quad B:X\to Z$$ be linear operators.

Then if $\ker A\subset\ker B,$ there exists an operator $\Lambda:Y\to Z$ such that $B=\Lambda A$:
\[
  \begin{tikzcd}
    X \arrow{r}{A} \arrow[swap]{dr}{B} & Y \arrow{d}{\Lambda} \\
     & Z
  \end{tikzcd}
\]
\end{lem}
From this and formulas (\ref{sdg-0}), (\ref{sfg43}), it follows 
that there is a row vector $\Lambda=(\lambda_1,\ldots,\lambda_n)(t,x,\dot x)$ such that 
\begin{equation}\label{dfg54hh}N=\Lambda\ph_{\dot x}.\end{equation}
Let us differentiate $\ph$ by virtue of system (\ref{azzxxx}) and equate the result to zero:
$$\ph_t+\ph_x\dot x+\ph_{\dot x}G^{-1}(f^T+N^T)=0.$$
Substituting $N$ from formula (\ref{dfg54hh}) here, we find
\begin{equation}\label{xdfgaxv}\Lambda^T=-\big(\ph_{\dot x}G^{-1}\ph_{\dot x}^T\big)^{-1}(\ph_t+\ph_x\dot x+\ph_{\dot x}G^{-1}f^T).\end{equation}
Here we use the following fact from linear algebra.
\begin{lem}\label{xfg6zzzz6}If a square $m\times m$ matrix $P$ is symmetric and positive definite, and matrix $B$ consists of $m$ columns and $n$ rows ($m>n$), with $\mathrm{rang}\,B=n,$ then the matrix $BPB^T$ is symmetric and positive definite.\end{lem}
The theorem is proved.

Equation (\ref{azzxxx}) with the substituted reaction (\ref {dfg54hh}) is called the Lagrange equation with multipliers or the Lagrange equation of the first kind.

\begin{theo}\label{dfg44}
Let the function $x(t)$ satisfy (\ref {xfg55}) and 
\begin{align}\ph_{\dot x}&\big(t,x(t),\dot x(t)\big)\xi=0\Longrightarrow \nonumber\\
&([\mathscr T]-f)\xi=0.\label{xsdfg23ee}\end{align}
Then $x(t)$ satisfies (\ref{azzxxx}) with reactions (\ref{dfg54hh}), (\ref{xdfgaxv}).
\end{theo}Equation (\ref{xsdfg23ee}) with virtual displacements $\xi$ is called the general equation of dynamics.

If one sequentially substitutes the basis vectors of the space of virtual displacements into (\ref{xsdfg23ee}), a system of $r$ scalar differential equations is obtained, each containing components of the vector $\ddot x$. Together with the constraint equations, this yields a system of $m$ differential equations of order $2r+n=2m-n=\dim S-1$.


{\it Proof of Theorem \ref{dfg44}.} Indeed, by Lemma \ref{dfh6ll} there exists a row vector
$\tilde\Lambda(t)$ such that 
\begin{equation}\label{dfh6pp}
[\mathscr T]-f=\ddot x^T(t)G-f\big(t,x(t),\dot x(t)\big)=\tilde\Lambda(t)\ph_{\dot x}\big(t,x(t),\dot x(t)\big).\end{equation}
Let us differentiate the equation $\ph\big(t,x(t),\dot x(t)\big)=0:$
$$\ph_t\big(t,x(t),\dot x(t)\big)+\ph_x\big(t,x(t),\dot x(t)\big)\dot x(t)+\ph_{\dot x}\big(t,x(t),\dot x(t)\big)\ddot x(t)=0.$$
Substituting $\ddot x$ from (\ref{dfh6pp}) here, we verify that 
$\tilde \Lambda(t)=\Lambda\big(t,x(t),\dot x(t)\big).$

The theorem is proved.




The following theorem is obvious.
\begin{theo}\label{dfg4xde4}
Let $x(t)$ satisfy (\ref{azzxxx}) with reactions (\ref{dfg54hh}), (\ref{xdfgaxv}).
Then the following implication holds: $$\ph_{\dot x}\big(t,x(t),\dot x(t)\big)\xi=0\Longrightarrow
([\mathscr T]-f)\xi=0.$$\end{theo}

Theorems \ref{as11}, \ref{dfg44}, \ref{dfg4xde4} constitute the D'Alembert-Lagrange principle.





\section{On the well-posedness of the definition of reactions and virtual displacements}\label{sdf55}
The definitions of reactions $N$ and the space of virtual displacements involve the functions $\ph$.
In this section, we show that neither the reactions nor the space of virtual displacements depends on the choice of functions $\ph$ defining the manifold $S$. In other words, reactions and virtual displacements depend on the geometry of the manifold $S$, and not on the specific method of its analytical specification.




Let us introduce a vector of smooth functions defined on $I\times D\times\mathbb{R}^m\times\mathbb{R}^n:$
$$U=(U^1,\ldots,U^n)^T(t,x,\dot x,z),\quad z\in\mathbb{R}^n,$$
such that
$$U(t,x,\dot x,z)=0\Leftrightarrow z=0,\quad \det U_z(t,x,\dot x,0)\ne 0,\quad \forall (t,x,\dot x)\in I\times D\times\mathbb{R}^m.$$
Let $\psi(t,x,\dot x):=U\big(t,x,\dot x,\ph(t,x,\dot x)\big).$ Then $S=\{\ph=0\}=\{\psi=0\}$
and the following formulas hold
$$\psi_t=U_t(t,x,\dot x,\ph)+U_z(t,x,\dot x,\ph)\ph_t,\quad U_t(t,x,\dot x,\ph)\mid_{\ph=0}=0.$$
Similarly
$\psi_x\mid_S=U_z\ph_x\mid_S,\quad \psi_{\dot x}\mid_S=U_z\ph_{\dot x}\mid_S.$

In particular $$\mathrm{rang}\,\big(\psi_{\dot x}\mid_S\big)=n, \quad \ker \big(\psi_{\dot x}\mid_S\big)=\ker\big( \ph_{\dot x}\mid_S\big).$$

Let us calculate the reactions corresponding to different ways of defining the manifold $S$:
\begin{align}
N^T&=-\ph^T_{\dot x}\big(\ph_{\dot x}G^{-1}\ph_{\dot x}^T\big)^{-1}(\ph_t+\ph_x\dot x+\ph_{\dot x}G^{-1}f^T)\nonumber\\
\tilde N^T&=-\psi^T_{\dot x}\big(\psi_{\dot x}G^{-1}\psi_{\dot x}^T\big)^{-1}(\psi_t+\psi_x\dot x+\psi_{\dot x}G^{-1}f^T).\nonumber
\end{align}
From this
$
N^T\mid_S=\tilde N^T\mid_S.$
\section{Holonomic and non-holonomic constraints}
\begin{df}Constraints (\ref{xfg55}) are called holonomic if there are functions
$$g=(g^1,\ldots, g^n)^T(t,x),\quad \mathrm{rang}\, g_x=n$$ such that 
$$S=\Big\{(t,x,\dot x)\mid\frac{d}{dt}g(t,x) =0\Big\},\quad \frac{d}{dt}g(t,x)=g_t(t,x)+g_x(t,x)\dot x.$$

Constraints that are not holonomic are called non-holonomic.
\end{df}
The question of the holonomicity of constraints is equivalent to the question of the integrability of a differential system on a manifold \cite{stern}.

The equations $g(t,x)=0$ are called geometric constraints, and the equations (\ref{xfg55}) are called differential constraints.

Hereafter, until the end of Section \ref{cbbbbb}, we will consider the constraints holonomic and write
\begin{equation}\label{xdfg11gy}\ph(t,x,\dot x)=g_t(t,x)+g_x(t,x)\dot x=0.\end{equation}
In particular,
\begin{equation}\label{xf11q}\ph_{\dot x}=g_x,\quad\ker \ph_{\dot x}(t,x,\dot x)=\ker g_x(t,x).\end{equation}



\section{Covariance of the variational derivative}
Assume that the manifold
$$\Sigma_t=\{x\in D\mid g(t,x)=0\},\quad \dim\Sigma_t=r$$ is the image of a smooth manifold $Y,\quad \dim Y=r$ under the embedding\footnote{By an embedding we mean a smooth injective mapping
$u(t,\cdot):Y\to\mathbb{R}^m$ such that for every point $y\in Y$ and its image $x=u(t,y)\in \mathbb{R}^m$ there exist neighborhoods $U_y\subset Y$ and $U_x\subset\mathbb{R}^m$ with local coordinates $(y^1,\ldots,y^r)$ and $(x^1,\ldots,x^m)$ respectively, in which $u$ takes the form: $x^i=y^i,\quad i=1,\ldots, r,\quad x^s=0,\quad s=r+1,\ldots,m,$ moreover $U_x\cap \Sigma_t=\{x^s=0,\quad s=r+1,\ldots,m\}.$
The specified local coordinates smoothly depend on $t\in I$.}
$$u(t,\cdot):Y\to D,\quad \Sigma_t=u(t,Y).$$
Let $y=(y^1,\ldots,y^r)^T$ denote the local coordinates in $Y$;
\begin{equation}\label{fg11}g(t,u(t,y))=0,\quad u=(u^1,\ldots,u^m)^T, \quad x=u(t,y).\end{equation}
The linear mapping
$$u_y(t,y):T_yY\to T_{u(t,y)}\Sigma_t$$ is an isomorphism, in particular,
$\mathrm{rang}\,u_y=r$,
and the tangent space
$T_{u(t,y)}\Sigma_t$ coincides with the space of virtual displacements:
$$T_{u(t,y)}\Sigma_t=\ker \big(\ph_{\dot x}\mid_{x=u(t,y)}\big).$$
This is clear if we differentiate formula (\ref{fg11}) with respect to $y$ and use (\ref{xf11q}):
$$g_x(t,u(t,y))u_y(t,y)=0.$$
Consider a smooth function $\mathscr{F}(t,x,\dot x)$ defined on the extended phase space.
\begin{df}
The variational derivative of $\mathscr{F}$ is defined as the following differential operator
$$[\mathscr{F}]=([\mathscr{F}]_{1},\ldots ,[\mathscr{F}]_{m}),\quad [\mathscr{F}]_{k}=\frac{d}{dt}\frac{\partial \mathscr{F}}{\partial\dot x^k}-\frac{\partial \mathscr{F}}{\partial x^k}.$$
\end{df}
Obviously, the variational derivative is a linear operation and
for any function $w=w(t,x),\quad w:I\times D\to \mathbb{R}$
the identity 
$$\Big[\frac{dw}{dt}\Big]=0$$
holds.



\begin{theo}\label{dfggg456}
The equality
$$[\mathscr{F}]\Big|_{x=u(t,y)}u_y(t,y)=[F]$$
holds, where the function $F:I\times TY\to\mathbb{R}$ ($TY$ is the tangent bundle) is defined by the formula
\begin{equation}\label{dff}F(t,y,\dot y)=\mathscr{F}\Big|_{x=u(t,y)}=\mathscr{F}\big(t,u(t,y),u_t(t,y)+u_y(t,y)\dot y\big).\end{equation}
\end{theo}


{\it Proof of Theorem \ref{dfggg456}.} Let us carry out the calculations in coordinates. By formula (\ref{dff}) we have:
$$\frac{\partial F}{\partial y^i}=\frac{\partial \mathscr{F}}{\partial x^p}\frac{\partial u^p}{\partial y^i}+
\frac{\partial \mathscr{F}}{\partial \dot x^p}\Big(\frac{\partial^2 u^p}{\partial y^l\partial y^i}\dot y^l+\frac{\partial^2u^p}{\partial t\partial y^i}\Big);$$
and
$$
\frac{d}{dt}\frac{\partial F}{\partial \dot y^i}=\frac{d}{dt}\Big(\frac{\partial \mathscr{F}}{\partial \dot x^p}
\frac{\partial u^p}{\partial y^i}\Big)=\frac{\partial u^p}{\partial y^i}\Big(\frac{d}{dt}\frac{\partial \mathscr{F}}{\partial \dot x^p}\Big)+\frac{\partial \mathscr{F}}{\partial \dot x^p}\Big(\frac{\partial^2 u^p}{\partial y^l\partial y^i}\dot y^l+\frac{\partial^2u^p}{\partial t\partial y^i}\Big).$$
Subtracting one equality from the other, we obtain the required result.
The theorem is proved.




\section{Lagrange equations}\label{cbbbbb}
Theorem \ref{dfggg456} means that the variational derivative behaves like a covector field with respect to the mapping $u$.
Roughly speaking, the derivation of the Lagrange equations consists of using the mapping $u$ to perform a pullback of the general equation of dynamics from the domain $D$ to the manifold $Y$.

Let us proceed to rigorous formulations.

\begin{df}Generalized forces are the components of the row vector
\begin{equation}\label{sfg40}Q(t,y,\dot y):=f\big(t,u(t,y),u_t(t,y)+u_y(t,y)\dot y\big) u_y(t,y).\end{equation}
\end{df}
\begin{theo}\label{xsfgtyuio}
Let the function $y(t)$ be a solution of the Lagrange equations (of the second kind)
\begin{equation}\label{sdfgqas}[T]=Q,\end{equation}
where the kinetic energy $T=T(t,y,\dot y)$ is defined by the formula:
\begin{equation}\label{xdfsgggtyu}
T=\mathscr T\mid_{x=u(t,y)}=
\frac{1}{2}\big(u_t+u_y\dot y\big)^TG\big(u_t+u_y\dot y\big).
\end{equation}

Then the function $x(t)=u(t,y(t))$ satisfies the general equation of dynamics (\ref{xsdfg23ee}) with constraints (\ref{xdfg11gy}).\end{theo}
{\it Proof of Theorem \ref{xsfgtyuio}.} By Theorem \ref{dfggg456}, we have
\begin{equation}\label{xdfg6yiii}([\mathscr{T}]-f)\Big|_{x=u(t,y(t))}u_y(t,y(t))=([T]-Q)\Big|_{y=y(t)}.\end{equation}
The right side of this equality is zero by the condition. As already noted, the image of the operator $u_y$ coincides with the space of virtual displacements, so the general equation of dynamics is satisfied.

The theorem is proved.

The manifold $Y$ is called the configuration manifold, and the local coordinates $y$ are called generalized coordinates of the system (\ref{sdfgqas}).

The manifolds $TY,\quad I\times TY$ are called the phase space and the extended phase space, respectively.

The components of the vector $\dot y$ are called generalized velocities.




\begin{theo}\label{dfb00dd}
Let the function $x(t)$ satisfy the general equation of dynamics (\ref{xsdfg23ee}) with constraints (\ref{xdfg11gy}), and at some $t_0\in I$ the equality
$g(t_0,x(t_0))=0$ holds.

Then there exists a unique solution $y(t)$ of equations (\ref{sdfgqas}) such that $$x(t)=u(t,y(t)).$$\end{theo}
{\it Proof of Theorem \ref{dfb00dd}}.
We integrate the equality $$\frac{d}{dt}g(t,x(t))=0$$ from $t_0$ to $t:\quad g(t,x(t))-g(t_0,x(t_0))=0.$
Thus, $x(t)\in \Sigma_t,\quad t\in I$.

For any $t$, the equation \begin{equation}\label{dh0009}x(t)=u(t,y)\end{equation} is uniquely solvable for $y$, since $u(t,\cdot):Y\to\Sigma_t$ is a diffeomorphism. Let us denote this solution as $y(t),\quad x(t)=u(t,y(t))$.

We will show that $y(t)$ is a smooth function. Indeed, (\ref{dh0009}) is a system of $m$ equations with an $r$-dimensional unknown vector $y$. We select $r$ independent equations from this system and apply the implicit function theorem regarding the smoothness of the solution with respect to the parameter.


The left side of equality (\ref{xdfg6yiii}) is zero by the condition.
The theorem is proved.



\begin{theo}\label{asdf4eew}The formula
$T=T_2+T_1+T_0$ holds, where $$T_2=\frac{1}{2}\dot y^Tu_y^TGu_y\dot y$$ is a positive definite quadratic form of the generalized velocities; $T_1=u_t^TGu_y\dot y$ is a linear form, $T_0=u_t^TGu_t/2$.\end{theo}
This statement follows directly from formula (\ref{xdfsgggtyu}) and Lemma \ref{xfg6zzzz6}.

From Theorem (\ref{asdf4eew}) it follows that system (\ref{sdfgqas}) can be represented in normal form:
$$\ddot y=a(t,y,\dot y),$$ and therefore the Cauchy existence and uniqueness theorem is valid for it.

If the forces $f$ are generalized potential forces, then the forces $Q$ are also generalized potential forces -- due to the covariance of the operation $[\cdot]$. Moreover
$Q=[V],\quad V(t,y,\dot y)=\mathscr V\mid_{x=u(t,y)},$ and equations (\ref{sdfgqas}) turn into Lagrange equations:
$$[L]=0,\quad L=T-V.$$The function $L=L(t,y,\dot y)$ is called the Lagrange function or Lagrangian.

\section{On classical notations}
Traditionally, in mechanics, the vector of virtual displacements $\xi$ is written as follows:
$$\xi=(\delta \boldsymbol r_1,\ldots,\delta\boldsymbol  r_\nu)^T.$$
Accordingly, equations (\ref{sdg-0}) and (\ref{sfg43}) have the form
$$\sum_{i=1}^\nu\Big(\frac{\partial \Phi^k}{\partial\boldsymbol{\dot r}_i},\delta \boldsymbol r_i\Big)=0,\quad\sum_{i=1}^\nu(\boldsymbol R_i,\delta \boldsymbol r_i)=0,\quad k=1,\ldots,n.$$



The general equation of dynamics (\ref{xsdfg23ee}) has the form
$$\sum_{i=1}^\nu(m_i\boldsymbol{\ddot r}_i-\boldsymbol F_i,\delta\boldsymbol r_i)=0.$$



The embedding $u$ is introduced as follows:
$$u(t,\cdot):y\mapsto (\boldsymbol r_1(t,y),\ldots,\boldsymbol r_\nu(t,y)).$$
Accordingly, the generalized forces (\ref{sfg40})
are calculated by the formula:
$$Q_j=\sum_{i=1}^\nu\Big(\bs F_i,\frac {\partial\bs r_i}{\partial  y^j}\Big),\quad Q=(Q_1,\ldots,Q_r).$$




\begin{thebibliography}{99}

\bibitem{stern} S. Sternberg: Lectures on Differential Geometry. Prentice-Hall,
Englewood  Cliffs, N. J., 1964.
 \end{thebibliography}

\end{document}
